\documentclass[11pt, a4paper]{article}

\usepackage[utf8]{inputenc}
\usepackage[T1]{fontenc}
\usepackage{amsmath, amssymb}
\usepackage{booktabs}
\usepackage{hyperref}
\usepackage{enumitem}
\usepackage{geometry}
\usepackage{natbib}
\usepackage{graphicx}
\usepackage{xcolor}

\geometry{margin=1in}
\hypersetup{colorlinks=true, linkcolor=blue, citecolor=blue, urlcolor=blue}

\title{\Large Emergent Duality: 14 Fundamental Dualities Across 6 Algebraic Layers,\\
Validated by Domain Analysis and Neurosymbolic Learning}

\author{
    J. Arturo Ornelas Brand \\
    \textit{Independent Researcher} \\
    \texttt{arturoornelas62@gmail.com}
}

\date{March 2026}

\begin{document}

\maketitle

\begin{abstract}
We present a formal ontological framework in which conceptual complexity emerges from 14 fundamental dualities organized across 6 algebraic layers (Boolean, Fuzzy, Ordinal, Modal, Trivalent, Probabilistic). The framework is grounded in 72 semantic primitives connected by a directed acyclic graph (DAG) of 133+ dependency edges. We validate the framework through three independent methodologies: (1)~domain analysis across 9 disciplines using the Information Domain Validity Score (IDVS), achieving IDVS~=~1.000 for all 8 positive domains against a null baseline ($p_{95}$~=~0.896, permutation $p$~=~0.0001); (2)~Bayesian model comparison yielding a Bayes factor of 121.4 (decisive) for the proposed ordering over reversed alternatives; and (3)~neurosymbolic learning via a GPT-2 Medium model augmented with a 72-bit triadic head, which learns the primitive structure at 90.8\% bit accuracy while preserving general language ability (PPL~=~31.95, identical to baseline). The learned representations exhibit near-perfect relational structure (Regla de Tres cosine~=~0.996) and show a significant phase transition during training ($p$~=~0.0004). We report both confirmatory and null results transparently: layer ordering in neural representations is not significant ($p$~=~0.71), and only 5 of 13 dual axes show significant learned anti-correlation after FDR correction.
\end{abstract}

\noindent\textbf{Keywords:} ontology, semantic primitives, emergence, duality, neurosymbolic AI, algebraic structure, domain validation

% ============================================================
\section{Introduction}

The search for universal semantic primitives---irreducible conceptual atoms from which all meaning is composed---spans philosophy (Leibniz's \emph{characteristica universalis}), linguistics (Wierzbicka's Natural Semantic Metalanguage), and computational semantics (ConceptNet, FrameNet). Existing approaches either lack formal algebraic structure or fail to demonstrate empirical universality across domains.

We propose that conceptual complexity emerges through a specific sequence of \textbf{dualities}---paired oppositions that each require all previous dualities as prerequisites. This emergence follows an algebraic chain:

$$\text{Boolean} \to \text{Fuzzy} \to \text{Ordinal} \to \text{Modal} \to \text{Trivalent} \to \text{Probabilistic}$$

Each algebra introduces operators that the previous one cannot express: degree from existence, comparison from degree, possibility from order, valuation from possibility, observation from valuation.

The contribution of this work is threefold:

\begin{enumerate}[nosep, leftmargin=*]
    \item \textbf{Formalization}: 72 semantic primitives organized in a DAG across 6 algebraic layers, with 14 dual axes and formal dependency rules.
    \item \textbf{Empirical validation}: cross-domain analysis (8 disciplines + 1 negative control) using a novel metric (IDVS) with statistical hardening (permutation tests, FDR correction, Bayesian model comparison, null baselines).
    \item \textbf{Neural validation}: a neurosymbolic model (GPT-2 Medium + triadic head) that learns the proposed structure from natural language, demonstrating that the algebraic relationships are recoverable from distributional semantics.
\end{enumerate}

% ============================================================
\section{The Framework}

\subsection{Semantic Primitives}

We define 72 semantic primitives distributed across 6 layers:

\begin{table}[ht]
\centering
\small
\begin{tabular}{llrl}
\toprule
\textbf{Layer} & \textbf{Algebra} & \textbf{Prims} & \textbf{Examples} \\
\midrule
1. Point (0D)    & Boolean $\{0,1\}$     & 3  & vac\'{i}o, informaci\'{o}n, uno \\
2. Line (1D)     & Fuzzy $\{0,+\}$       & 8  & fuerza, contenci\'{o}n, uni\'{o}n, separaci\'{o}n \\
3. Time (1D+t)   & Ordinal $\{<,=,>\}$   & 13 & mover, orden, caos, creaci\'{o}n \\
4. Plane (2D)    & Modal $\{\Diamond,\Box\}$ & 21 & bien, mal, verdad, mentira \\
5. Volume (3D)   & Trivalent $\{-,0,+\}$ & 23 & vida, muerte, placer, dolor \\
6. Observer (3D+) & Probabilistic $\{0,?\}$ & 4 & temporal\_obs, eterno\_obs \\
\bottomrule
\end{tabular}
\caption{72 semantic primitives across 6 algebraic layers.}
\label{tab:primitives}
\end{table}

Each primitive is assigned a unique prime number, a bit position, a layer, a set of dependencies (parents in the DAG), and a natural-language definition.

\subsection{The Dependency DAG}

Primitives are connected by a DAG with 133+ directed edges encoding the relation ``X requires Y to exist.'' Key structural properties:

\begin{itemize}[nosep, leftmargin=*]
    \item \textbf{Layer monotonicity}: all dependencies of a layer-$n$ primitive lie in layers $\leq n$.
    \item \textbf{3 roots}: vac\'{i}o (bit 0), informaci\'{o}n (bit 1), uno (bit 44)---all layer~1.
    \item \textbf{Transitive closure}: the maximum ancestor set size is 22 (for creador\_obs, layer~6).
\end{itemize}

\subsection{The 14 Dualities}

The 14 dualities emerge in a specific order, each anchored by 2--4 primitives:

\begin{table}[ht]
\centering
\small
\begin{tabular}{clll}
\toprule
\textbf{\#} & \textbf{Duality} & \textbf{Anchors} & \textbf{Layer} \\
\midrule
D1  & Existir / No-existir     & vac\'{i}o, informaci\'{o}n, uno           & 1 \\
D2  & Aqu\'{i} / No-aqu\'{i}  & eje\_profundidad, eje\_vertical            & 2--4 \\
D3  & Antes / Despu\'{e}s     & posici\'{o}n\_temporal, flujo\_temporal    & 3 \\
D4  & Posible / Imposible      & puede, debe, tal\_vez                     & 4 \\
D5  & Identidad / Diferencia   & muchos, tipo\_de, todo, algunos           & 2--4 \\
D6  & Movimiento / Quietud     & mover, hacer, quietud                     & 3 \\
D7  & Orden / Caos             & orden, caos, bien, mal, verdad, mentira   & 3--4 \\
D8  & Uno / Muchos             & uno, muchos                               & 1--2 \\
D9  & Dentro / Fuera           & contenci\'{o}n, uni\'{o}n, separaci\'{o}n & 2 \\
D10 & Parte / Todo             & parte\_de, todo, algunos                  & 2--4 \\
D11 & Vida / Muerte            & vida, muerte, placer, dolor               & 5 \\
D12 & Causa / Efecto           & porque, si\_entonces                      & 3 \\
D13 & Semejante / Diferente    & tipo\_de                                  & 4 \\
D14 & Sujeto / Objeto          & consciente, ausente, temporal\_obs         & 5--6 \\
\bottomrule
\end{tabular}
\caption{The 14 fundamental dualities with anchor primitives and layer assignments.}
\label{tab:dualities}
\end{table}

\subsection{Algebraic Layer Chain}

Each layer operates with a distinct algebra. The key insight is that each transition is \textbf{irreversible}: you cannot define the operators of layer~$n$ without having the operators of layer~$n{-}1$. This is verified computationally at 86\% match (12/14 dualities correctly predicted by their algebra assignment).

% ============================================================
\section{Domain Validation}

\subsection{The IDVS Metric}

We define the \textbf{Information Domain Validity Score} for a domain~$d$ as:

$$\text{IDVS}(d) = \text{coverage}(d) \times \text{monotonicity}(d)$$

\noindent where:
\begin{itemize}[nosep, leftmargin=*]
    \item \textbf{Coverage} = $1 - \frac{\text{nulls in layers 1--4}}{\text{total primitives in layers 1--4}}$. A domain that classifies foundational primitives as N (``does not apply'') is penalized.
    \item \textbf{Monotonicity} = $1 - \frac{\text{violations}}{\text{total edges}}$. A violation occurs when a child primitive has a higher state than its parent ($D > A > N$).
\end{itemize}

\subsection{Classification Protocol}

Each of 72 primitives is classified in each domain as:
\begin{itemize}[nosep, leftmargin=*]
    \item \textbf{D} (Direct): the primitive operates directly in this domain.
    \item \textbf{A} (Analogical): the primitive applies by metaphor or analogy.
    \item \textbf{N} (Null): the primitive does not apply.
\end{itemize}

\subsection{Results}

All 8 positive domains achieve IDVS~=~1.000 after D-A-N reclassification:

\begin{table}[ht]
\centering
\small
\begin{tabular}{llcccc}
\toprule
\textbf{Domain} & \textbf{Type} & \textbf{IDVS} & \textbf{Cov.} & \textbf{Mono.} & \textbf{Perm.~$p$} \\
\midrule
Music       & Positive  & 1.000 & 1.000 & 1.000 & 0.0001 \\
Chemistry   & Positive  & 1.000 & 1.000 & 1.000 & 0.0001 \\
Biology     & Positive  & 1.000 & 1.000 & 1.000 & 0.0001 \\
Mathematics & Positive  & 1.000 & 1.000 & 1.000 & 0.0001 \\
Philosophy  & Positive  & 1.000 & 1.000 & 1.000 & 0.0001 \\
Physics     & Positive  & 1.000 & 1.000 & 1.000 & 0.0001 \\
Linguistics & Positive  & 1.000 & 1.000 & 1.000 & 0.0001 \\
Psychology  & Positive  & 1.000 & 1.000 & 1.000 & 0.0001 \\
\textbf{Astrology} & \textbf{Negative} & \textbf{0.479} & 0.650 & 0.737 & 0.0001 \\
\bottomrule
\end{tabular}
\caption{IDVS results across 9 domains. The negative control (Astrology) correctly fails.}
\label{tab:idvs}
\end{table}

\subsection{Null Baseline}

To calibrate the IDVS threshold, we computed 10,000 random D-A-N assignments preserving the marginal distribution of each domain. The null distribution yields: mean max IDVS~=~0.861, $p_{95}$~=~0.896 (threshold set at 0.90), $p_{99}$~=~0.955. All 8 positive domains exceed $p_{99}$.

\subsection{Sensitivity Analysis}

\begin{table}[ht]
\centering
\small
\begin{tabular}{llrl}
\toprule
\textbf{Perturbation} & \textbf{Target} & \textbf{Mean $\Delta$IDVS} & \textbf{Large disruptions} \\
\midrule
Remove edge  & DAG topology        & $-0.003$ & 3.4\% \\
Move layer   & Layer assignment    & $0.000$  & 0.0\% \\
Flip D-A-N   & Classification logic & $-0.060$ & 71.4\% \\
\bottomrule
\end{tabular}
\caption{Sensitivity analysis. Robustness gradient: Topology $>$ Layering $>$ Logic.}
\label{tab:sensitivity}
\end{table}

\subsection{Bayesian Model Comparison}

We compared 4 orderings of the 14 dualities using triangularity of the dependency matrix as the test statistic:

\begin{table}[ht]
\centering
\small
\begin{tabular}{lcc}
\toprule
\textbf{Model} & \textbf{Triangularity} & \textbf{BF vs.\ Reversed} \\
\midrule
\textbf{Proposed} & \textbf{0.539} & \textbf{121.4} (decisive) \\
Hegel    & 0.528 & --- \\
Peirce   & 0.484 & --- \\
Reversed & 0.143 & 1.0 \\
\bottomrule
\end{tabular}
\caption{Bayesian model comparison of duality orderings.}
\label{tab:bayesian}
\end{table}

Against 1,000 random orderings ($\mu = 0.338$), the proposed ordering is significant ($p = 0.005$).

% ============================================================
\section{Neural Validation}

\subsection{Architecture}

We augment GPT-2 Medium (355M parameters) with a \textbf{triadic head}: a single linear projection from the 1024-dimensional hidden state to a 72-bit semantic vector. The activation function is iFSQ ($2\sigma(1.6x) - 1$), which distributes activations uniformly across $[-1, +1]$.

\noindent\textbf{Total trainable parameters (v5\_frozen)}: 73,728 (0.02\% of model). The GPT-2 backbone is frozen.

\subsection{Training}

\begin{table}[ht]
\centering
\small
\begin{tabular}{ll}
\toprule
\textbf{Component} & \textbf{Configuration} \\
\midrule
Loss       & $\mathcal{L}_\text{lang} + 0.05 \cdot (\mathcal{L}_\text{tri} + 2\mathcal{L}_\text{sup} + 5\mathcal{L}_\text{sub})$ \\
$\mathcal{L}_\text{tri}$  & diversity + InfoNCE + entropy + alignment \\
$\mathcal{L}_\text{sup}$  & MSE(projection, gold\_target) \\
$\mathcal{L}_\text{sub}$  & $\text{relu}(\text{hypernym}_{01} - \text{hyponym}_{01}).\text{mean}()$ \\
Optimizer  & AdamW ($3 \times 10^{-4}$, $\beta = 0.9/0.95$) \\
Schedule   & Cosine decay, 50\% triadic warmup \\
Data       & WikiText-103 ($\sim$100M tokens) \\
Steps      & 50,000 \\
Hardware   & NVIDIA RTX 4060 Ti 16GB \\
\bottomrule
\end{tabular}
\caption{Training configuration for v5\_frozen.}
\label{tab:training}
\end{table}

\subsection{Training Run Summary}

\begin{table}[ht]
\centering
\small
\begin{tabular}{lcrrcrl}
\toprule
\textbf{Run} & \textbf{Bits} & \textbf{Acc.} & \textbf{Unique} & \textbf{Coh.} & \textbf{PPL} & \textbf{Key Change} \\
\midrule
v1 & 65 & $\sim$90\% & ? & ? & 2,777 & Baseline \\
v2 & 65 & 95.8\% & 3.8\% & ? & 6,735 & Bad targets \\
v3 & 65 & 90.9\% & 89.2\% & 0.653 & 5,388 & Fixed: 65 prims \\
v4 & 72 & 92.5\% & 100\% & 0.876 & 5,608 & +7 primitives \\
\textbf{v5} & \textbf{72} & \textbf{90.8\%} & \textbf{100\%} & \textbf{0.818} & \textbf{31.95} & \textbf{Frozen base} \\
\bottomrule
\end{tabular}
\caption{Training run progression. V5 solves catastrophic forgetting via frozen backbone.}
\label{tab:runs}
\end{table}

\subsection{Key Finding: Field Density $>$ Individual Targets}

V4 added 7 new primitives to differentiate collapsed dual pairs (Jaccard $> 0.80$). The result was surprising: collapsed pairs resolved even though their individual gold targets did not change. The mechanism is \textbf{field density}---more diverse training examples teach the model what good differentiation looks like.

\subsection{Key Finding: Catastrophic Forgetting Solved}

All runs v1--v4 showed catastrophic language degradation (PPL $> 2{,}700$ vs.\ baseline 31.95). V5 froze the GPT-2 backbone, training only the 73K-parameter triadic head. Result: PPL~=~31.95 (identical to baseline) at a cost of only $-1.7\%$ accuracy.

\subsection{Relational Structure (Regla de Tres)}

The model preserves proportional relationships between dual pairs:

\begin{table}[ht]
\centering
\small
\begin{tabular}{lcc}
\toprule
\textbf{Analogy} & \textbf{Cosine} & \textbf{Bit Acc.} \\
\midrule
good:evil = truth:lie                    & 0.998 & 1.000 \\
creation:destruction = life:death        & 0.998 & 1.000 \\
freedom:control = individual:collective  & 0.995 & 0.986 \\
pleasure:pain = conscious:absent         & 0.996 & 1.000 \\
union:separation = order:chaos           & 0.997 & 1.000 \\
creation:destruction = union:separation  & 0.997 & 1.000 \\
\bottomrule
\end{tabular}
\caption{Regla de Tres: relational structure preservation. Mean cosine: 0.996.}
\label{tab:r3}
\end{table}

\subsection{Neural Probes}

8 falsifiable questions were tested using reptimeline checkpoint analysis:

\begin{table}[ht]
\centering
\small
\begin{tabular}{clll}
\toprule
\textbf{Probe} & \textbf{Question} & \textbf{Result} & \textbf{$p$-value} \\
\midrule
Q1 & Layer order in representations?   & \textbf{Null}              & 0.713 \\
Q2 & Dual anti-correlation?            & \textbf{5/13 significant}  & $<0.018$ (BH) \\
Q3 & DAG recovery?                     & \textbf{Above random}      & 0.001 \\
Q5 & Phase transition?                 & \textbf{Yes, step 6000}    & 0.0004 \\
Q4 & Triadic dependencies?             & Null                       & 1.000 \\
Q6 & Unsupervised discovery?           & Null                       & 1.000 \\
Q7 & Causal verification?              & Null                       & 1.000 \\
\bottomrule
\end{tabular}
\caption{Neural probe results. The model learns relational but not ordinal structure.}
\label{tab:probes}
\end{table}

\noindent\textbf{Interpretation}: The model learns \emph{relational} structure (analogies, some dual axes, DAG edges) but does not recover \emph{ordinal} structure (layer ordering, ternary dependencies). This is consistent with the frozen-base design: the head projects semantics into bit space without restructuring the backbone's internal representations.

% ============================================================
\section{Discussion}

\subsection{What the Framework Claims}

The Emergent Duality framework makes two testable claims:

\begin{enumerate}[nosep, leftmargin=*]
    \item \textbf{Structural claim}: 72 primitives connected by a specific DAG can classify concepts across any knowledge domain using a ternary (D/A/N) code.
    \item \textbf{Ordering claim}: the 14 dualities emerge in a specific sequence dictated by algebraic dependency.
\end{enumerate}

\subsection{Evidence For}

\begin{itemize}[nosep, leftmargin=*]
    \item \textbf{Structural claim}: supported. 8/8 positive domains achieve IDVS~=~1.000, all exceeding the $p_{99}$ null threshold (0.955). The negative control correctly fails.
    \item \textbf{Ordering claim}: partially supported. Bayesian comparison yields BF~=~121.4 (decisive) against reversed ordering, but only weak evidence against Hegel (BF~=~1.14) and Peirce (BF~=~1.95). The neural model does not recover layer ordering (Q1, $p$~=~0.71).
\end{itemize}

\subsection{Honest Limitations}

\begin{enumerate}[nosep, leftmargin=*]
    \item \textbf{Self-evaluated classifications}: All D-A-N assignments were made by the framework author. Inter-rater data exists for only 1/9 domains (Mathematics). The permutation tests validate internal consistency but not external agreement.
    \item \textbf{Underpowered rank correlations}: With $n = 7$ (circle) or $n = 14$ (combined), Spearman correlations cannot achieve significance at $\alpha = 0.05$ even for moderate effects.
    \item \textbf{IDVS~=~1.000 for all positives}: After reclassification (v1.1), all positive domains achieve perfect scores. This may reflect overfitting to the metric rather than genuine universality.
    \item \textbf{Neural probes}: 5/8 probes returned null results. The model learns surface relational structure but not deep algebraic properties.
    \item \textbf{Scale--algebra tradeoff}: Larger models (355M) show lower subsumption rates (76.9\%) than smaller ones (40M, 98.3\%).
\end{enumerate}

\subsection{Future Directions}

This work opens several lines of investigation:

\begin{enumerate}[nosep, leftmargin=*]
    \item \textbf{Independent rater validation}: The most critical next step. Blind inter-rater studies across all 9 domains would transform self-evaluated D-A-N classifications into externally validated ones. Preliminary data for Mathematics (1 rater) shows agreement; 8 domains remain.
    \item \textbf{Scale--algebra tradeoff}: Understanding why 355M-parameter models show lower subsumption (76.9\%) than 40M models (98.3\%) may reveal fundamental constraints on discrete algebraic learning at scale.
    \item \textbf{Cross-lingual validation}: All primitives are currently named and evaluated in a Spanish-language conceptual frame. Testing whether the same DAG structure holds when primitives are derived independently from other linguistic traditions would strengthen universality claims.
    \item \textbf{Relational prime chains}: Extending the triadic head with compositional prime chain verification for O(1) anti-hallucination---a natural application of the algebraic structure to production AI systems.
    \item \textbf{Expanded corpus training}: V5 was trained on WikiText-103. Training on multilingual or domain-specific corpora could reveal whether the 72-primitive structure generalizes beyond English-language distributional semantics.
\end{enumerate}

% ============================================================
\section{Related Work}

\begin{itemize}[nosep, leftmargin=*]
    \item \textbf{Wierzbicka (1996)}: Natural Semantic Metalanguage proposes $\sim$65 universal primitives but lacks formal algebraic structure or computational validation.
    \item \textbf{ConceptNet} \citep{speer2017conceptnet}: Large-scale knowledge graph with 34 relation types but no emergence hierarchy.
    \item \textbf{FrameNet} \citep{baker1998framenet}: Frame-based semantics covering $\sim$1,200 frames but domain-specific rather than universal.
    \item \textbf{Category theory approaches} \citep{coecke2010mathematical}: DisCoCat and related formalisms provide compositional semantics but do not address emergence ordering.
\end{itemize}

Our framework differs in proposing (a)~a specific algebraic chain governing emergence, (b)~empirical validation across 9 domains, and (c)~neurosymbolic learnability from natural language.

% ============================================================
\section{Conclusion}

We have presented a formal ontological framework of 14 fundamental dualities emerging across 6 algebraic layers, grounded in 72 semantic primitives. The framework achieves perfect IDVS scores across 8 disciplines while correctly rejecting a negative control, with statistical hardening via permutation tests, null baselines, and Bayesian model comparison. A neurosymbolic model (GPT-2 + triadic head) learns the proposed structure at 90.8\% accuracy without degrading language ability, exhibiting near-perfect relational preservation (cosine 0.996).

The main open question is external validation: independent human evaluators must confirm the D-A-N classifications before universality claims can be considered robust. We release all data, scripts, and model code to enable replication.

\subsection*{Software and Data Availability}

All code and data are available under BUSL-1.1 at:
\begin{itemize}[nosep, leftmargin=*]
    \item Validation scripts, DAG data, and domain analyses: \texttt{triadic-emergent-duality} repository.
    \item Neurosymbolic training pipeline (GPT-2 + triadic head): \texttt{triadic-microgpt} repository.
    \item Representation lifecycle tracking: \texttt{reptimeline} (also available on PyPI).
    \item Triadic projection head (standalone): \texttt{triadic-head} (also available on PyPI).
\end{itemize}
The repository also includes a structured compilation of first principles across 14 scientific domains (92 files organized hierarchically from mathematics to interdisciplinary fields), used by the algebraic gap analysis to search for candidate operators in uncovered inter-layer connections.

% ============================================================
\appendix
\section{Notation}

\begin{table}[ht]
\centering
\small
\begin{tabular}{ll}
\toprule
\textbf{Symbol} & \textbf{Meaning} \\
\midrule
D, A, N & Direct, Analogical, Null classification states \\
IDVS    & Information Domain Validity Score \\
BF      & Bayes Factor \\
PPL     & Perplexity \\
iFSQ    & Improved Finite Scalar Quantization \\
BH-FDR  & Benjamini--Hochberg False Discovery Rate \\
SHD     & Structural Hamming Distance \\
DAG     & Directed Acyclic Graph \\
\bottomrule
\end{tabular}
\caption{Notation reference.}
\label{tab:notation}
\end{table}

% ============================================================
\bibliographystyle{plainnat}

\begin{thebibliography}{7}

\bibitem[Wierzbicka(1996)]{wierzbicka1996}
A.~Wierzbicka.
\newblock \emph{Semantics: Primes and Universals}.
\newblock Oxford University Press, 1996.

\bibitem[Speer et~al.(2017)]{speer2017conceptnet}
R.~Speer, J.~Chin, and C.~Havasi.
\newblock ConceptNet 5.5: An open multilingual graph of general knowledge.
\newblock In \emph{AAAI}, 2017.

\bibitem[Baker et~al.(1998)]{baker1998framenet}
C.~F. Baker, C.~J. Fillmore, and J.~B. Lowe.
\newblock The Berkeley FrameNet project.
\newblock In \emph{COLING-ACL}, 1998.

\bibitem[Coecke et~al.(2010)]{coecke2010mathematical}
B.~Coecke, M.~Sadrzadeh, and S.~Clark.
\newblock Mathematical foundations for a compositional distributional model of meaning.
\newblock \emph{Linguistic Analysis}, 2010.

\bibitem[Radford et~al.(2019)]{radford2019language}
A.~Radford, J.~Wu, R.~Child, et~al.
\newblock Language models are unsupervised multitask learners.
\newblock \emph{OpenAI Technical Report}, 2019.

\bibitem[Benjamini and Hochberg(1995)]{benjamini1995controlling}
Y.~Benjamini and Y.~Hochberg.
\newblock Controlling the false discovery rate: A practical and powerful approach to multiple testing.
\newblock \emph{Journal of the Royal Statistical Society: Series B}, 57(1):289--300, 1995.

\bibitem[Mentzer et~al.(2025)]{mentzer2025fsq}
F.~Mentzer et~al.
\newblock Finite scalar quantization: {VQ-VAE} made simple.
\newblock In \emph{ICLR}, 2024.

\end{thebibliography}

\end{document}
